Grading risk, view accuracy, and weak-target overlap assess severity prediction, acquisition-view recognition, and agreement with local annotations, respectively. The full model's lowest grading risk and the deletion variants' higher spatial-task scores motivate evaluating the score and each accompanying output together. The module-deletion comparisons evaluate whole designs rather than isolating supervision, capacity, or information use, leaving the mechanisms behind these trade-offs unresolved.

Local rectangular boxes provide positional cues for spatial learning. Comparing regional maps with the original image and clinician boxes makes correspondence to these cues visible and identifies additional regions for inspection. Weak-target overlap measures positional agreement, whereas complete lesion coverage is a separate objective. Responses outside local boxes may reflect relevant tissue or erroneous activation; more complete spatial references and clinician review could distinguish these possibilities.

For clinical review, the joint outputs provide three objects to inspect: the predicted view, regional response, and final score. The illustrated view error despite close score agreement shows how these outputs expose different aspects of a prediction. Displaying them alongside the original image would allow separate checks of anatomical context and local responses. Their practical benefit remains to be evaluated through reader studies of error identification, appropriate assessment revisions, and review time.

Deployment studies should measure preprocessing, inference, and display together on target hardware, beyond the model-only costs reported here. Independent-cohort evaluation should test grading and spatial outputs across different patient populations and acquisition settings.
